% =============================================================================
\chapter{Interface e Experiência do Usuário (UI/UX) Científica}

\section{Design Centrado em Dados}

A interface do \texttt{app.py} organiza o fluxo analítico de forma hierárquica,
priorizando a clareza e a produtividade do cientista em detrimento de elementos
decorativos. A distribuição dos elementos entre abas, colunas e área central segue
princípios de design centrado no usuário, com decisões arquiteturais explícitas sobre
onde cada tipo de conteúdo deve residir.

\begin{description}

\item[a)] \textbf{Navegação por Etapas (\texttt{st.tabs})}: As três abas —
\textbf{Configuração}, \textbf{Resultados} e \textbf{Gráficos} — guiam o usuário em um
fluxo linear e irreversível: entrada de dados $\rightarrow$ execução $\rightarrow$
inspeção de resultados. Essa estrutura espelha o \textit{workflow} típico de validação de
modelos quimiométricos e reduz a sobrecarga cognitiva ao esconder informações irrelevantes
para a etapa atual.

\item[b)] \textbf{Layout em Colunas e Hierarquia Visual}: O \texttt{app.py} utiliza
\texttt{st.columns(2)} para agrupar uploads de calibração e teste lado a lado,
sinalizando visualmente a paridade entre os conjuntos. A área central é reservada para
cards de métricas gerados por \texttt{components.metric\_card\_html()} e para a tabela
comparativa em \texttt{pd.DataFrame}. O CSS centralizado em \texttt{components.APP\_CSS}
— incluindo gradientes de fundo, bordas coloridas por modelo e estilo terminal para o
log — replica dashboards técnicos e industriais sem poluir o \texttt{app.py} com
definições de estilo.

\item[c)] \textbf{Feedback Progressivo}: O botão \textbf{Executar Análise} é
desabilitado (\texttt{disabled=not files\_ok}) enquanto algum dos quatro arquivos estiver
ausente, prevenindo execuções inválidas. A barra de progresso atualizada pelo callback e
a mensagem \texttt{st.success} ao término fornecem confirmação clara de conclusão,
padrão recomendado pelas diretrizes de usabilidade para interfaces científicas.

\item[d)] \textbf{Decisão Arquitetural: \texttt{st.tabs} vs.\ \texttt{st.sidebar}}: O
Streamlit oferece dois paradigmas de organização de parâmetros técnicos e resultados:
a \texttt{st.sidebar} (painel lateral persistente, visível em todas as telas) e o
\texttt{st.tabs} (abas sequenciais que segmentam o fluxo). Para o perfil de usuário
científico deste sistema — que opera em etapas ordenadas e não retorna aos parâmetros
após a execução — o modelo de abas foi preferido por três razões: (1) impõe
explicitamente a sequência \textit{Configuração $\rightarrow$ Resultados $\rightarrow$
Gráficos}, prevenindo execuções prematuras; (2) preserva a área central integralmente
para visualização de resultados e gráficos, sem que uma barra lateral compita com os
gráficos do Plotly em largura de tela; e (3) alinha-se ao padrão de \textit{wizards}
adotado em \textit{softwares} de análise científica (Origin, MATLAB App Designer), nos
quais cada etapa revela seu conteúdo apenas quando os pré-requisitos anteriores foram
cumpridos. A \texttt{st.sidebar} permanece como alternativa recomendada para versões
futuras que implementem parâmetros globais persistentes, como configuração de idioma,
normalização de dados ou seleção de paleta de cores para publicação.

\end{description}

% ─────────────────────────────────────────────────────────────────────────────
\section{Visualização Interativa com Plotly}

A escolha do Plotly no \texttt{components.py} é fundamentada pela necessidade de
inspeção técnica precisa em espectrometria e regressão multimodelo.

\begin{description}

\item[a)] \textbf{Integração Nativa e Centralizada}: Com a separação em
\texttt{components.py}, todas as configurações visuais — cores por modelo
(\texttt{COR\_RLM}, \texttt{COR\_BAGGING}, \texttt{COR\_SUBAGGING}), estilos de eixo
(\texttt{\_AXIS\_STYLE}), cores de fundo (\texttt{\_PLOT\_BG}) — estão definidas em um
único arquivo. Alterar a paleta de cores ou o estilo de grade para padrões de publicação
científica requer modificação em um só lugar, sem tocar no \texttt{app.py} ou no
\texttt{models.py}.

\item[b)] \textbf{Ferramentas de Inspeção Científica}: Os gráficos \texttt{go.Scatter}
gerados por \texttt{build\_chart()} e \texttt{build\_scatter\_chart()} habilitam zoom,
\textit{pan}, \textit{hover} com coordenadas exatas e legendas interativas. A linha ideal
\textit{y=x} (\texttt{mode="lines", dash="dot"}) no gráfico de dispersão
(\texttt{build\_scatter\_chart}) permite validação visual imediata de viés e dispersão
dos três modelos sobrepostos — RLM Simples, Bagging e Subagging — em uma única figura.

\item[c)] \textbf{Padrões de Publicação}: Os eixos customizados (\texttt{showgrid=True},
\texttt{gridcolor="\#000000"}, fontes pretas explícitas via \texttt{\_AXIS\_STYLE})
garantem legibilidade para relatórios técnicos, apresentações e submissão a periódicos,
alinhando-se às normas de visualização de dados científicos sem exigir pós-processamento
externo das figuras.

\end{description}

% ─────────────────────────────────────────────────────────────────────────────
\section{Contratos de Interface e Funções}
\label{sec:contratos}

A separação em três módulos aplica rigorosamente o princípio de separação de
responsabilidades, tornando o sistema testável, manutenível e extensível
independentemente em cada camada.

\begin{description}

\item[a)] \textbf{Separação de Camadas}: O \texttt{models.py} não importa
\texttt{streamlit} nem \texttt{plotly}. O \texttt{components.py} não importa
\texttt{streamlit} nem \texttt{numpy}. O \texttt{app.py} não contém nenhuma função
matemática — a instrução mais complexa é a chamada \texttt{result =
run\_full\_analysis(...)}. Essa separação permite testar o \texttt{models.py}
isoladamente com \texttt{pytest} e dados sintéticos, sem instanciar nenhum componente de
interface.

\item[b)] \textbf{Contratos Explícitos via \textit{Type Hints}}: Todas as funções do
\texttt{models.py} declaram tipos de entrada e saída:
\texttt{perform\_bagging(...) -> np.ndarray},
\texttt{calc\_metrics(...) -> tuple[float, float, float]},
\texttt{find\_optimal\_parameters(...) -> tuple[dict, str]}.
No \texttt{app.py}, \texttt{parse\_int\_list()} e \texttt{parse\_float\_list()} lançam
\texttt{ValueError} com mensagens orientativas para entradas fora do domínio válido —
inteiros positivos para \textit{m} e decimais em $(0,1]$ para \textit{k} — prevenindo
falhas silenciosas na camada de negócio.

\item[c)] \textbf{Manutenibilidade e Escalabilidade}: Adicionar um novo modelo preditivo
(como Ridge Regression ou Random Forest) requer apenas uma nova função em
\texttt{models.py} e um novo bloco de exibição em \texttt{app.py}. Modificar a paleta de
cores ou o layout dos gráficos requer apenas alterações em \texttt{components.py}.
Nenhuma dessas mudanças impacta as outras camadas, confirmando que a arquitetura modular
adotada atende ao critério de \textbf{abertura para extensão, fechamento para
modificação} (\textit{Open/Closed Principle}) proposto por Robert C.
Martin~\cite{martin2008}.

\end{description}
